% !TeX root = ../main.tex
\section{Signal processing primer}
\label{sec:dsp}

The pieces of a guitar rig fall into four kinds, and each kind calls for a
different estimation strategy. This section explains them.

\subsection{Linear filters: EQ, tone stack, cabinet}

A processor is \term{linear and time-invariant} (LTI) when (a) doubling the
input doubles the output, (b) the response to a sum of inputs is the sum of the
responses, and (c) it behaves the same at all times. Equalisers, the tone stack
of an amplifier (taken alone), a speaker cabinet with its microphone, and to a
good approximation a room, are LTI.

\begin{keyidea}[: an LTI system is fully described by its impulse response]
Feed an LTI system a single click (an \term{impulse}); the output is its
\term{impulse response} (IR) $h$. The output for any input $x$ is then the
\term{convolution}
\[
  y[n] = \sum_{k} h[k]\,x[n-k],
\]
each input sample triggering a scaled copy of the IR. Equivalently, in
frequency, each frequency is multiplied by the \term{frequency response}
$H(f)$, the Fourier transform of $h$. This is why a cabinet can be stored as a
\file{.wav} IR and why the app's Cabinet pedal is a convolution
(\file{backend/dsp/convolution.rs}).
\end{keyidea}

Two consequences matter for tone estimation.
\begin{itemize}
\item \textbf{The linear part can be measured from long-term spectra.} If
  $y = h * x$, the average spectrum of $y$ is the average spectrum of $x$
  times $|H(f)|^2$, so in decibels
  $L_y(f) = L_x(f) + 20\log_{10}|H(f)|$. The difference between two
  \term{long-term average spectra} is the EQ that turns one into the other.
  This is what Tone Match products and the app's final step do.
\item \textbf{Phase.} Many IRs have the same magnitude response $|H(f)|$. The
  \term{minimum-phase} one concentrates its energy as early as possible, so it
  adds the least delay and smearing~\cite{oppenheim2010dtsp}. When only a
  magnitude is known (as from spectra), the minimum-phase IR is the natural
  choice. It can be computed through the \term{cepstrum} (the inverse Fourier
  transform of the log magnitude): keeping only its causal half and
  transforming back gives the minimum-phase filter (Section~\ref{sec:impl-match}).
\end{itemize}

\subsection{Nonlinear processors: drive, distortion, the amplifier}

A distortion circuit \term{clips}: above some level, the output stops growing.
Clipping a sine wave adds harmonics (multiples of its frequency); clipping two
notes together adds sum and difference frequencies (\term{intermodulation}).
Two properties of nonlinear processors drive the design of the app's search.

\begin{keyidea}[: input level is a setting]
How much a clipper distorts depends on how loud its input is. The same amp at
the same gain sounds clean with a soft guitar and crunchy with a hot pickup. A
tone estimate must therefore choose the input level together with the gain; in
the app it is the first searched parameter (\code{input level (RMS)}).
\end{keyidea}

\begin{worked}[: order matters]
Put a treble boost \emph{before} a distortion: the high strings drive the
clipper harder, so they distort more and gain harmonics, and the result is
brighter and fizzier. Put the same boost \emph{after}: the distortion is the
same as without it and only the final colour is brighter. For an LTI chain, the
order of filters does not matter; around a nonlinearity it does. Corollary: a
linear filter \emph{after} all the clipping (the cabinet, mix EQ, and in the
app's amp model the tone stack) can be estimated separately at the end, which is
what the app does.
\end{worked}

\subsection{Dynamics: compressors}

A \term{compressor} lowers the gain when the signal exceeds a
\term{threshold}, by a \term{ratio} (4:1 means 4\dB{} above the threshold
becomes 1\dB), reacting with an \term{attack} time and recovering with a
\term{release} time. It changes the envelope of notes (less attack, more
sustain) more than their spectrum. It is nonlinear and time-varying, and it is
easy to confuse with the natural compression of a distorted amp.

\subsection{Time-varying effects: modulation}

Chorus, flanger, phaser and tremolo are driven by a \term{low-frequency
oscillator} (LFO, typically 0.1--10\Hz) that slowly moves a delay time, a
filter notch or the volume. Their sound depends on the LFO rate and depth. The
search landscape over an LFO rate is very rugged (many local minima): a small
change in rate shifts the LFO's phase completely after a few seconds.

\subsection{Long-memory effects: delay and reverb}

A \term{delay} adds copies of the signal after a fixed time (the \term{delay
time}), each copy quieter by the \term{feedback} factor. A \term{reverb} adds a
dense, decaying tail; its decay time (RT60, the time to fall by 60\dB) and
brightness are its main traits. Like LFO rates, delay times make rugged
landscapes: a delay of 300\ms{} and one of 310\ms{} produce echoes that land in
different places, so the comparison score does not change smoothly between
them. That is why the app \emph{measures} the delay time from the recording
instead of searching it (Section~\ref{sec:impl-delay}).

\subsection{Loudness}

Two renders that differ only in volume should count as the same tone, because
the mix level of a guitar says nothing about its rig. All comparisons in the app
first scale both signals to the same RMS (root-mean-square) level.
